\section{AR 模型用历史预测未来}
\label{sec:06-Random-Processes/06-ARMA-and-ARIMA/01}

假设你拿到一串按月记录的数据——某种商品的月销量，过去五年的 60 个数字。扫一眼就有感觉：上个月卖得多，这个月大概也不差；连续三个月上涨，下个月未必还能涨那么多。

你的直觉在说一件事：历史值携带了关于未来的信息。但"携带多少"和"怎样携带"这两问，直觉回答不了。

自回归（AR）模型是处理这个问题的最简框架。它的核心假设只有一句话：下一个观测值可以写成最近几个历史值的线性组合，再加上一个无法用历史解释的随机创新。

\subsection{最简单的尝试：用昨天预测今天}

从 AR(1) 开始——只用一个历史值：

\[
X_t=\phi X_{t-1}+\varepsilon_t,
\]

其中 \(\varepsilon_t\) 是白噪声：零均值、方差 \(\sigma^2\)、不同时刻不相关。\(\phi\) 是一个落在 \((-1,1)\) 里的实数。

如果 \(\phi=0.8\)，今天大致是昨天的 80\%，加上一点新鲜扰动。如果 \(\phi=-0.5\)，今天的符号大致和昨天相反——序列在两个方向之间摆动。如果 \(\phi=1.02\)，序列会爆炸；今天比昨天多走的那 2\%，明天会叠加上去，后天的基数又变大——系统不会稳定在某个均值附近。

参数 \(\phi\) 于是承载了两个含义：它度量了记忆强度（越接近 \(\pm1\)，记忆越长），也决定了过程是否平稳。

从 \(X_t=\phi X_{t-1}+\varepsilon_t\) 反复回代，可以把 \(X_t\) 写成过去所有创新的加权和：

\[
X_t=\varepsilon_t+\phi\varepsilon_{t-1}+\phi^2\varepsilon_{t-2}+\cdots
=\sum_{j=0}^{\infty}\phi^j\varepsilon_{t-j}.
\]

这个展开只有当 \(|\phi|<1\) 时才收敛（权重平方可和），过程才有有限方差。此时方差的解析表达式是 \(\Var(X_t)=\sigma^2/(1-\phi^2)\)，自相关函数呈指数衰减：\(\rho_k=\phi^{|k|}\)。\(\phi\to1\) 时方差发散——随机游走卡在平稳世界的边界上。

\subsection{不只昨天：AR(p) 与特征多项式}

一个滞后不够怎么办？把昨天、前天、大前天都拉进来：

\[
X_t=\phi_1X_{t-1}+\phi_2X_{t-2}+\cdots+\phi_pX_{t-p}+\varepsilon_t.
\]

定义滞后算子 \(B\)（\(BX_t=X_{t-1}\)，\(B^2X_t=X_{t-2}\)，依此类推），AR(p) 可以压缩成一句：

\[
\phi(B)X_t=\varepsilon_t,\qquad
\phi(z)=1-\phi_1z-\phi_2z^2-\cdots-\phi_pz^p.
\]

\(\phi(z)\) 叫特征多项式。平稳性条件推广了 AR(1) 的 \(|\phi|<1\) 到了高维：\(\phi(z)=0\) 的所有根必须满足 \(|z|>1\)，即全在复平面的单位圆外。

为什么要关心这个条件？写出 AR(p) 的因果表示——把 \(X_t\) 表为白噪声的单边滑动平均：

\[
X_t=\frac{1}{\phi(B)}\varepsilon_t
=\psi(B)\varepsilon_t
=\sum_{j=0}^{\infty}\psi_j\varepsilon_{t-j}.
\]

\(\psi(B)=1/\phi(B)\) 要能展开成收敛的幂级数，就要求 \(\phi(z)\neq0\) 对 \(|z|\le1\) 中的所有 \(z\) 成立。根在单位圆外等价于所有极点落在单位圆内——从复分析的角度，这保证了有理函数的幂级数展开在单位圆盘上解析，从而 \(\psi_j\) 绝对可和。

用 AR(2) 做具体检查。特征方程 \(1-\phi_1z-\phi_2z^2=0\) 的根在单位圆外这件事，等价于三个实数不等式：

\[
\phi_2+\phi_1<1,\qquad
\phi_2-\phi_1<1,\qquad
|\phi_2|<1.
\]

这三个条件是 \((\phi_1,\phi_2)\) 平面上的一个倒三角区域。参数落在这个三角里，过程平稳；落在外面，过程或发散或表现出其他非平稳行为。

\subsection{ACF 与 PACF：两个工具，一个目的——找 \(p\)}

知道了模型形式，下一个问题是：数据应该塞进 AR(几)？

自相关函数（ACF）\(\rho_k=\gamma_k/\gamma_0\) 衡量 \(X_t\) 和 \(X_{t-k}\) 之间的相关性——这是两个时刻的直接相关，不扣掉中间任何东西。AR(p) 的 ACF 不会在某个滞后突然归零，而是按照差分方程的解衰减：对实根是指数衰减，对复根是振荡衰减。衰减速度由最靠近单位圆的那个根决定。

但 ACF 会骗人。假设真实的 \(X_t=0.8X_{t-1}+\varepsilon_t\)（纯 AR(1)），那么 \(X_t\) 和 \(X_{t-2}\) 也会表现相关——因为 \(X_t\) 和 \(X_{t-1}\) 相关、\(X_{t-1}\) 又和 \(X_{t-2}\) 相关。这个两步传导会产生一个非零的 \(\rho_2\)，即便 \(X_{t-2}\) 对 \(X_t\) 没有直接作用。

偏自相关函数（PACF）\(\alpha_k\) 就是来解决"传导 vs 直接"这个混淆的。它回答的是：在已经控制了中间所有滞后 \(X_{t-1},\ldots,X_{t-k+1}\) 的前提下，\(X_{t-k}\) 对 \(X_t\) 还有多少剩余相关性？翻译成人话：去掉了所有间接传话之后，第 \(k\) 步的历史还有没有独家的新信息。

PACF 有一个干净的性质：对 AR(p)，\(\alpha_k=0\) 对所有 \(k>p\) 成立——截尾。解释很简单：AR(p) 的性质决定了，控制了中间 \(p\) 步以后，更早的历史在给定这 \(p\) 步的条件下与当前值是条件独立的。PACF 的截尾点就直接告诉你阶数。

实际计算 PACF 的方法是连续拟合 AR(1), AR(2), AR(3), \ldots，每次取最后一个系数的估计 \(\widehat\phi_{kk}\) 作为 \(\alpha_k\)。样本 PACF 在真阶之后的滞后上应围绕零随机波动，近似标准差 \(1/\sqrt{n}\)。如果 \(k=3\) 时 PACF 还远远超出 \(\pm1.96/\sqrt{n}\) 的区间，说明 \(p=2\) 不够。

\subsection{Yule--Walker 方程：架起 ACF 与参数之间的桥}

ACF 告诉你数据长什么样，Yule--Walker 方程告诉你怎样从 ACF 倒推出参数。推导步骤只需两步：用 \(X_{t-k}\) 乘 AR(p) 方程两边，取期望。

对于 \(k\ge1\)，白噪声项 \(\varepsilon_t\) 与 \(X_{t-k}\) 不相关（因为因果表示里 \(X_{t-k}\) 只包含 \(\varepsilon_{t-k}\) 及更早的创新），所以 \(\E[\varepsilon_t X_{t-k}]=0\)。剩下的是自协方差的递推关系：

\[
\gamma_k=\phi_1\gamma_{k-1}+\phi_2\gamma_{k-2}+\cdots+\phi_p\gamma_{k-p},\qquad k\ge1.
\]

取 \(k=1,\ldots,p\)，这组方程可以用矩阵写出：

\[
\begin{pmatrix}
\gamma_1\\ \gamma_2\\ \vdots\\ \gamma_p
\end{pmatrix}
=
\begin{pmatrix}
\gamma_0 & \gamma_1 & \cdots & \gamma_{p-1}\\
\gamma_1 & \gamma_0 & \cdots & \gamma_{p-2}\\
\vdots & \vdots & \ddots & \vdots\\
\gamma_{p-1} & \gamma_{p-2} & \cdots & \gamma_0
\end{pmatrix}
\begin{pmatrix}
\phi_1\\ \phi_2\\ \vdots\\ \phi_p
\end{pmatrix}.
\]

左边的协方差矩阵是对称 Toeplitz 结构——每条反对角线上的元素相同。给定样本自协方差 \(\widehat\gamma_k\)，解这个线性系统就得到 Yule--Walker 估计 \(\widehat\phi_j\)。

Toeplitz 矩阵是正定的，所以解在理论上总是存在且唯一。而且在足够大的样本下，由正定样本 Toeplitz 矩阵给出的 Yule--Walker 解倾向于落在平稳区域内——这是它相较于普通最小二乘的一个实用优势。渐近分布上，Yule--Walker 与条件最小二乘收敛到同一个正态极限，但有限样本下二者可以不相等；选择哪个取决于你更信任少偏还是更低均方误差。

\subsection{一个可以手算的 AR(2)}

取 \(\rho_1=0.8,\;\rho_2=0.65\)。AR(2) 的 Yule--Walker 方程为

\[
\begin{aligned}
\rho_1&=\phi_1+\phi_2\rho_1,\\
\rho_2&=\phi_1\rho_1+\phi_2.
\end{aligned}
\]

第一个方程给出 \(\phi_1=\rho_1(1-\phi_2)\)。代入第二个：

\[
\rho_2=\rho_1(1-\phi_2)\rho_1+\phi_2
=\rho_1^2(1-\phi_2)+\phi_2
=\rho_1^2+(\phi_2-\rho_1^2\phi_2).
\]

整理：

\[
\phi_2=\frac{\rho_2-\rho_1^2}{1-\rho_1^2}
=\frac{0.65-0.64}{1-0.64}
=\frac{0.01}{0.36}\approx0.028.
\]

回代得 \(\phi_1=0.8\times(1-0.028)\approx0.778\)。

系数 \(\phi_2\) 非常小——第二个滞后几乎没有额外的直接贡献。这和直觉一致：\(\rho_1=0.8,\rho_2=0.65\)，而 AR(1) 的 \(\rho_2\) 理论值是 \(\rho_1^2=0.64\)；观测到的 0.65 只比 0.64 多出 0.01，正是 \(\phi_2=0.028\) 传导过去的那一点额外作用。这个例子把 PACF 为什么能在滞后 2 附近归零这件事，用具体数字而不是抽象概念展示了出来。

当然，并非任意一对 \((\rho_1,\rho_2)\) 都合法。自协方差序列必须对应一个正半定的 Toeplitz 矩阵——你无中生有编一组自相关系数，可能根本没有实数 AR 过程能产生它。

\subsection{AR 模型的边界：什么时候它不够用}

AR 好用，但有清楚的前提。下面四种情况会突破它的假设：

\begin{itemize}
\tightlist
\item
  \textbf{长记忆}：ACF 不按指数衰减，而是按幂律缓慢衰减。金融波动率的自相关有时用 AR(1) 拟合，残差却在大时间尺度上仍有结构。这种情况需要 ARFIMA（自回归分数整合滑动平均），用分数差分把记忆参数从整数解放到实数。
\item
  \textbf{非线性依赖}：条件均值不是历史的线性函数。比如"连续跌三天"和"连续涨三天"引发的后续行为不对称——门槛自回归（TAR）让参数随历史值所处区间变化，或者直接用神经网络学习条件均值。
\item
  \textbf{异方差}：白噪声的方差本身在变。AR 假设 \(\varepsilon_t\) 方差恒定，但金融数据常有波动率聚集——大波动之后还是大波动。GARCH 族模型把方差本身也建模为过去冲击的函数，Part 11 会详细展开。
\item
  \textbf{结构突变}：\(\phi_1,\ldots,\phi_p\) 随时间变化。政策突变、经济危机、传感器更换——都会让过去拟合的参数在现在失效。状态空间模型或变点检测方法把参数视为可随时间演化的隐藏状态。
\end{itemize}

AR 是时间序列分析的第一个驿站，不是终点站。它的力量在于：一旦确定序列平稳，AR 就是一个可解释、可诊断、可快速部署的基线。下一篇把滑动平均（MA）加进来——不是让历史直接影响今天，而是让过去的冲击在今天留下回声。
